\UseRawInputEncoding
\documentclass[runningheads]{llncs}
\usepackage[utf8]{inputenc}
\usepackage[portuguese]{babel}
\usepackage{geometry}
\usepackage{graphicx}
\usepackage{amsmath}
\usepackage{url}
\usepackage{booktabs}
\usepackage{listings}
\usepackage{xcolor}
\usepackage{float}


% Configuração para mostrar código Java
\definecolor{javared}{rgb}{0.6,0,0} % for strings
\definecolor{javagreen}{rgb}{0.25,0.5,0.35} % comments
\definecolor{javapurple}{rgb}{0.5,0,0.35} % keywords
\definecolor{javadocblue}{rgb}{0.25,0.35,0.75} % javadoc


\lstset{
    language=Java,
    basicstyle=\ttfamily\footnotesize,
    keywordstyle=\color{javapurple}\bfseries,
    stringstyle=\color{javared},
    commentstyle=\color{javagreen},
    morecomment=[s][\color{javadocblue}]{/**}{*/},
    numbers=left,
    numberstyle=\tiny\color{black},
    stepnumber=1,
    numbersep=10pt,
    tabsize=2,
    showspaces=false,
    showstringspaces=false,
    frame=single,
    breaklines=true,
    captionpos=b,
}


\begin{document}

\title{Sistemas Distribuídos - Trabalho Prático}

\author{Daniel Pereira\inst{106912} \and
Fernando Ferreira\inst{106878} \and
José Diogo Fernandes\inst{104}}

\authorrunning{Grupo }

\institute{Universidade do Minho, Braga, Portugal\\}

\maketitle

\begin{abstract}
Este relatório descreve a implementação de um sistema de serviço de registo de eventos em séries temporais e de agregação de informação, relativos a venda de produtos. 
  A informação é mantida num \emph{Servidor} e acedida remotamente através de um \emph{Cliente} conectado por sockets TCP.

  \keywords{TCP \and Serialização TLV \and Java Concurrency.}
\end{abstract}

\section{Introdução}
O sistema desenvolvido baseia-se numa arquitetura cliente-servidor, onde a comunicação é estabelecida através de sockets TCP. O servidor foi desenhado de forma a processar pedidos de multiplos clientes concurrentemente, pelo que foi necessário cuidados para garantir a sincronização de dados dentro das várias estruturas do servidor.
Esta solução proposta visa uma correção eficiente, valorizando estratégias que minimizam o número de threads ativas. 
\texttt{ALGUEM MELHORE ISTO AQUI PURFA.}

\section{Pedidos}
Os pedidos dos clientes ao servidor são realizados através do protocolo TCP, garantindo fiávelmente que o pedido e a sua resposta sejam recebidas.

Na interface do cliente todos os pedidos são \emph{case-insensitive} de forma a facilitar a escrita e apelar a todos os utilizadores.
Um cliente realiza um pedido através da interface dada, estes pedidos podem ser dos tipos:

\subsection{Autenticação}
Para poder realizar pedidos um utilizador tem, obrigatoriamente, de estar autenticado. Para este fim são disponibilizados os pedidos de \texttt{Register} e \texttt{Login}.
Ambos os pedidos necessitam de dois campos para serem válidos.

\begin{itemize}
  \item \emph{Username}: O nome pelo qual o utilizador será autenticado.
  \item \emph{Password}: A palavra-passe do utilizador.
\end{itemize}

Caso já autenticado um utilizador não poderá se autenticar novamente nem registar nenhum utilizador.


\subsection{Registo de Eventos}
O registo de um evento relativo ao dia corrente é feito através do pedido \texttt{ADD\_SALE}.
Este pedido será executado corretamente apenas quando o produto ao qual se pretende adicionar um evento existe, pelo que caso contrário nenhuma alteração será feita. Para este efeito são necessários os argumentos.

\begin{itemize}
  \item \emph{ProductName}: O nome do produto ao qual se pretende adicionar um evento.
  \item \emph{Quantity}: A quantidade de produtos vendida. 
  \item \emph{SellPrice}: O preço pelo qual os produtos foram vendidos.
\end{itemize}

\subsection{Queries}
Um cliente pode realizar \emph{Queries} sobre os dados do servidor.
Todas as \emph{queries} disponibilizadas necessitam dos mesmos dados para serem executadas corretamente:
\begin{itemize}
  \item \emph{ProductName}: O nome do produto sobre o qual se realiza a query.
  \item \emph{NumDays}: O número de dias posteriores ao atual sobre o qual se realiza a query.
\end{itemize}

\noindent Um cliente pode fazer pedidos relativos à quantidade de que um produto foi vendido, ao preço máximo pelo qual foi vendido, o preço médio e ainda a receita total das vendas desse produto.
Para isto são usados respetivamente os pedidos:

\noindent \texttt{Query\_QTD}
\\\texttt{Query\_MAX}
\\\texttt{Query\_AVG}
\\\texttt{Query\_TOTAL}

\subsection{Filtros}
Os pedidos de filtros são pedidos que dado um conjunto de itens e um dia anterior ao corrente, devolve uma lista de todos os eventos desses itens de forma organizada e compacta.
\begin{itemize}
  \item \emph{ProductNames}: Nomes dos produtos
  \item \emph{Day}: Dia sobre o qual se quer os eventos
\end{itemize}

\texttt{METER EXEMPLO DE COMO OS DADOS SAO RETORNADOS AQUI}

\subsection{Notificações}
\texttt{AINDA NAO IMPLEMENTADO!}

\end{document}
